%%%%%%%%%%%%%%%%%%%%%%%%%%%%%%%%%%%%%%%%%%%%%%%%%%%%%%%%%%%%%%%%%%%%%%%%%%%%%%%%
% Obxectivo: Lista de siglas, abreviaturas, acrónimos, etc. empregados         %
%            no documento, xunto cos seus respectivos significados.            %
%                                                                              %
% Formato: \newacronym{clave}{sigla}{forma desenvolvida}                       %
% Nota: glossaries só imprime as entradas referenciadas no texto (\gls{clave}, %
%       \acrshort{}, \acrlong{}...). Para forzar a aparición de todas as       %
%       entradas nesta versión preliminar, pódese engadir \glsaddall antes de  %
%       \printglossary en memoria_tfg.tex.                                     %
%                                                                              %
% Organización: entradas agrupadas por capítulo de primera aparición en        %
%               el texto. Las entradas sin referencia van al final.            %
%%%%%%%%%%%%%%%%%%%%%%%%%%%%%%%%%%%%%%%%%%%%%%%%%%%%%%%%%%%%%%%%%%%%%%%%%%%%%%%%

% -----------------------------------------------------------------------
% Cap. 1 — Introducción
% -----------------------------------------------------------------------

\newacronym{htn}{HTN}{\textit{Hierarchical Task Network} (Red de Tareas Jerárquicas)}
\newacronym{ia}{IA}{Inteligencia Artificial}
\newacronym{il}{IL}{\textit{Imitation Learning} (Aprendizaje por Imitación)}
\newacronym{rl}{RL}{\textit{Reinforcement Learning} (Aprendizaje por Refuerzo)}

% -----------------------------------------------------------------------
% Cap. 2 — Estado del Arte
% -----------------------------------------------------------------------

\newacronym{dagger}{DAgger}{\textit{Dataset Aggregation}}
\newacronym{dqn}{DQN}{\textit{Deep Q-Network}}
\newacronym{drl}{DRL}{\textit{Deep Reinforcement Learning} (Aprendizaje por Refuerzo Profundo)}
\newacronym{llm}{LLM}{\textit{Large Language Model} (Modelo de Lenguaje Grande)}
\newacronym{mas}{MAS}{\textit{Multi-Agent System} (Sistema Multiagente)}
\newacronym{vpt}{VPT}{\textit{Video PreTraining}}

% -----------------------------------------------------------------------
% Cap. 3 — Marco Teórico
% -----------------------------------------------------------------------

\newacronym{bc}{BC}{\textit{Behavioral Cloning}}
\newacronym{mdp}{MDP}{\textit{Markov Decision Process} (Proceso de Decisión de Markov)}
\newacronym{ppo}{PPO}{\textit{Proximal Policy Optimization}}
\newacronym{sac}{SAC}{\textit{Soft Actor-Critic}}
\newacronym{sat}{SAT}{\textit{Boolean Satisfiability Problem} (Problema de Satisfacibilidad Booleana)}
\newacronym{strips}{STRIPS}{\textit{Stanford Research Institute Problem Solver}}

% -----------------------------------------------------------------------
% Cap. 4 — Metodología y planificación
% -----------------------------------------------------------------------

\newacronym{http}{HTTP}{\textit{Hypertext Transfer Protocol} (Protocolo de Transferencia de Hipertexto)}
\newacronym{yolo}{YOLO}{\textit{You Only Look Once}}
\newacronym{cpu}{CPU}{\textit{Central Processing Unit} (Unidad Central de Procesamiento)}
\newacronym{gpu}{GPU}{\textit{Graphics Processing Unit} (Unidad de Procesamiento Gráfico)}
\newacronym{ram}{RAM}{\textit{Random Access Memory} (Memoria de Acceso Aleatorio)}

% -----------------------------------------------------------------------
% Cap. 5 — Arquitectura del Sistema
% -----------------------------------------------------------------------

\newacronym{api}{API}{\textit{Application Programming Interface} (Interfaz de Programación de Aplicaciones)}
\newacronym{fov}{FOV}{\textit{Field of View} (Campo de Visión)}
\newacronym{hud}{HUD}{Interfaz de Información en Pantalla, \textit{Heads-Up Display}}
\newacronym{json}{JSON}{\textit{JavaScript Object Notation}}
\newacronym{rgb}{RGB}{\textit{Red, Green, Blue} (Rojo, Verde, Azul)}

% -----------------------------------------------------------------------
% Cap. 7 — Iteración II - Agente de Imitación
% -----------------------------------------------------------------------

\newacronym{cnn}{CNN}{\textit{Convolutional Neural Network} (Red Neuronal Convolucional)}
\newacronym{convlstm}{ConvLSTM}{\textit{Convolutional Long Short-Term Memory}}
\newacronym{gru}{GRU}{\textit{Gated Recurrent Unit}}
\newacronym{rnn}{RNN}{\textit{Recurrent Neural Network} (Red Neuronal Recurrente)}
\newacronym{vit}{ViT}{\textit{Vision Transformer}}

% -----------------------------------------------------------------------
% Cap. 8 — Iteración III - Agente de Refuerzo
% -----------------------------------------------------------------------

\newacronym{gae}{GAE}{\textit{Generalized Advantage Estimation}}
